\title{DeepSeekMath: Pushing the Limits of Mathematical Reasoning in Open Language Models}
The filtering process operates as follows: if any segment of text within a page contains a 10-gram that precisely matches a substring found in the benchmark datasets, the page is removed from the math training corpus. We provide empirical evidence supporting this view, specifically in the context of mathematical reasoning: training with code improves model performance in mathematical problem-solving, both with and without auxiliary tool support. To reduce redundancy, URL-based deduplication and near-deduplication strategies are applied to the raw Common Crawl data, producing 40B HTML documents. Moreover, due to model size constraints, DeepSeekMath does not match GPT-4’s effectiveness in few-shot learning scenarios: GPT-4 exhibits enhanced performance with additional exemplars, whereas DeepSeekMath's zero-shot and few-shot outcomes remain largely similar. Contrary to common expectations, our empirical findings suggest that exposure to arXiv papers does not enhance mathematical reasoning ability. The results indicate that DeepSeekMath-Base achieves notable performance in few-shot autoformalization. Here, we present an RL approach termed Group Relative Policy Optimization (GRPO), designed to be both effective and resource-efficient. This pattern is expected—early in training, the actor closely resembles the SFT model, leading to only subtle differences in sampled data. For each domain, we assess what fraction of its web pages were retrieved during the first data collection iteration. \item \textbf{Multilingual}:
    DeepSeekMath~Corpus contains material spanning several languages, with English and Chinese as its primary components. Our results are compared with contemporaneous state-of-the-art models:

\begin{itemize}[topsep=0pt]
    \item \textbf{Closed-source models} evaluated include: (1) models from the GPT suite, especially GPT-4 \citep{gpt4} and GPT-4 Code Interpreter~\footnote{\url{https://openai.com/blog/chatgpt-plugins\#\#code-interpreter}}, (2) Gemini Ultra and Pro \citep{gemini}, (3) Inflection-2 \citep{inflection-2}, (4) Grok-1~\footnote{\url{https://x.ai/model-card}}, along with (5) Baichuan-3~\footnote{\url{https://www.baichuan-ai.com}}, and (6) the most recent GLM-4~\footnote{\url{https://open.bigmodel.cn/dev/api\#glm-4}} from the GLM family \citep{glm}, all of which have been newly released by various Chinese technology companies. \paragraph{Mathematical Problem Solving with Tool Use}  
We assess capabilities in program-based mathematical reasoning using GSM8K and MATH, adopting few-shot program-of-thought prompting methods \citep{pot,pal}. The entirety of these experiments utilizes the HAI-LLM framework \citep{haillm}, chosen for its efficiency and lightweight nature. We outline several future research avenues pertaining to each of these components. Table \ref{tab:understanding_reasoning_code} illustrates that DeepSeekMath-Base 7B achieves considerable improvements over its predecessor, DeepSeek-Coder-Base-v1.5 \citep{deepseek-coder}, particularly on the MMLU and BBH benchmarks, reflecting enhanced language understanding and reasoning due to math-focused training. As illustrated in Figure \ref{fig:rl-analysis}, GRPO consistently outperforms online RFT, thereby demonstrating the effectiveness of modifying positive and negative gradient coefficients. A fastText model \citep{joulin2016fasttext} is trained to identify additional web pages resembling those in OpenWebMath. Our evaluation framework encompasses an extensive analysis of DeepSeekMath-Base 7B’s mathematical competencies, emphasizing its proficiency in independently constructing mathematical solutions, leveraging external tools for problem solving, and performing formal theorem verification. Its optimization targets the maximization of the following surrogate loss:

\begin{equation}
\footnotesize
                \mathcal{J}_{PPO}(\theta) = \mathbb{E}{[q \sim P(Q), o \sim \pi_{\theta_{old}}(O|q)]} \frac{1}{|o|} \sum_{t=1}^{|o|} \min \left[ \frac{\pi_\theta(o_{t} | q, o_{<t})}{\pi_{\theta_{old}}(o_{t} | q, o_{<t})} A_{t}, \text{clip} \left( \frac{\pi_\theta(o_{t} | q, o_{<t})}{\pi_{\theta_{old}}(o_{t} | q, o_{<t})}, 1 - \varepsilon, 1 + \varepsilon \right)  A_{t} \right] ,
\end{equation}

Here, $\pi_{\theta}$ and $\pi_{\theta_{old}}$ denote the current and previous iterations of the policy model, respectively. To probe whether joint training on a blend of code and math tokens in a single phase simultaneously augments mathematical reasoning and counteracts catastrophic forgetting, we employ a mixed-token regime. \end{itemize}

These models are generally intended for a wide range of applications, with most having been subjected to various alignment strategies. Examining Equations \ref{eq:GC-RFT} and \ref{eq:GC-GRPO}, a crucial distinction emerges: GRPO tailors its gradient scaling dynamically in accordance with reward values from the reward model, thus enabling nuanced adjustment—more pronounced encouragement or discouragement—depending on the degree of response quality. 
\begin{abstract}
The intricacies and structured patterns inherent in mathematical reasoning present a formidable obstacle for language models. The findings show that pre-training on mathematical data enhances the model's proficiency in both language understanding and reasoning tasks. Our findings demonstrate GRPO’s utility, even when DeepSeekMath-Instruct 7B achieves strong benchmark results. This produces the set $\mathbf{R} = \{ \{r_1^{index(1)},\cdots,r_1^{index(K_1)}\}, \cdots,  \{r_G^{index(1)},\cdots,r_G^{index(K_G)}\} \}$, with $index(j)$ marking the ending token of step $j$ and $K_i$ being the total step count for $o_i$. RFT and DPO employ an offline sampling paradigm, whereas Online RFT and GRPO utilize online sampling. Next, we perform manual annotation to label the URLs within these math-related domains that correspond to mathematical content (e.g., \url{mathoverflow.net/questions}). \subsection{Perspectives on Reinforcement Learning}
\subsubsection{Towards a Cohesive Training Framework} Here, we aim to articulate a general framework for analyzing disparate training strategies including SFT, RFT, DPO, PPO, GRPO, and empirically investigate the factors that comprise this unification. Introducing code pretraining significantly benefits program-based mathematical reasoning in both sequential and joint training scenarios. \item Through our experiments, we show that GRPO markedly improves the performance of our instruction-tuned system, DeepSeekMath-Instruct, relying solely on instruction-tuning data. \item \textbf{Open-source models} span: foundational models such as (1) DeepSeek-LLM-Chat 67B \citep{deepseek-llm}, (2) Qwen 72B \citep{qwen}, (3) SeaLLM-v2 7B \citep{seallm}, and (4) ChatGLM3 6B \citep{chatglm3}; along with models featuring mathematical improvements, including (5) InternLM2-Math 20B~\footnote{\url{https://github.com/InternLM/InternLM-Math}}, derived from InternLM2 and subjected to mathematics-focused pre-training and subsequent instruction tuning, (6) Math-Shepherd-Mistral 7B, which incorporates PPO training \citep{schulman2017proximal} into Mistral 7B \citep{mistral} using a reward model guided by process supervision, (7) the WizardMath suite \citep{wizardmath}, which augments mathematical reasoning for both Mistral 7B and Llama-2 70B \citep{llama2} through evolve-instruct—a variant of instruction tuning that utilizes AI-evolved prompts—combined with PPO optimization on primarily GSM8K and MATH datasets, (8) MetaMath 70B \citep{metamath}, where Llama-2 70B is refined using an expanded dataset from GSM8K and MATH, (9) ToRA 34B \cite{tora}, a CodeLlama 34B model optimized for tool-augmented mathematical reasoning, and (10) MAmmoTH 70B \citep{MathInstruct}, a Llama-2 70B variant instruction-tuned with MathInstruct. Remarkably, on the challenging competition-grade MATH dataset, this model outperforms every other open-source system and the vast majority of closed-source alternatives (such as Inflection-2 and Gemini Pro) by at least 9\% in absolute terms. Through the development of a careful data filtration system, we assemble the DeepSeekMath Corpus, consisting of 120B high-quality tokens sourced from mathematically relevant web content. Notably, when code and math tokens are provided together in one-stage training, there is a marked reduction in catastrophic forgetting associated with staged approaches, and both code generation abilities (Table \ref{tab:code-math-general-eval}) and program-assisted mathematical problem-solving (Table \ref{tab:code-math}) are further enhanced. The maximum sequence length is fixed at 1024, with a batch size of 1024. For English, we utilize GSM8K \citep{gsm8k}, MATH \citep{MATH}, SAT \citep{llemma}, OCW Courses \citep{minerva}, and MMLU-STEM \citep{mmlu}; Chinese tasks include MGSM-zh \citep{mgsm}, CMATH \citep{wei2023cmath}, Gaokao-MathCloze \citep{agieval}, and Gaokao-MathQA \citep{agieval}. \end{itemize}

\subsection{Summary of Evaluations and Metrics}

\begin{itemize}[topsep=0pt]
    \item \textbf{English and Chinese Mathematical Reasoning}: Our models are rigorously evaluated on a diverse suite of benchmarks in both English and Chinese, addressing mathematical challenges that span from elementary school to university curricula. \label{eq:objective}
\end{equation}

This formulation consists of three primary elements: 1) the \textit{Data Source} $\mathcal{D}$, responsible for specifying the set of examples used during training; 2) the \textit{Reward Function} $\pi_{{rf}}$, which supplies the evaluative feedback driving the learning process; 3) the \textit{Algorithm} $\mathcal{A}$, which utilizes both the data and the reward information to compute the gradient coefficient $GC$ that controls the strength of updates (either penalizing or reinforcing). While DeepSeekMath delivers strong results on quantitative reasoning assessments, its performance is notably less robust in geometric problem-solving and theorem-proof domains, especially when compared to closed-source models. For instance, the PRM800K datasets \citep{lightman2023let}, even though meticulously annotated by expert annotators, still exhibit around 20\% annotation errors\footnote{\url{https://github.com/openai/prm800k/issues/12\#issuecomment-1728491852}}. We initialize the reward model using DeepSeekMath-Base 7B with a 2e-5 learning rate. These candidates are then ranked according to the fastText model's scoring, and only the highest-ranking entries are retained to ensure quality. These results not only outperform all open-source models within the 7B–70B parameter range, but also exceed those of most closed-source systems. Additionally, by continuously incorporating code-related tokens during further training, DeepSeekMath-Base 7B preserves the high-level coding performance characteristic of DeepSeek-Coder-Base-v1.5 on both programming tasks. \item \textbf{Chinese mathematical datasets}: Our Chinese dataset contains K-12 mathematics problems spanning 76 subfields, including topics like linear equations. With further mathematical instruction tuning and reinforcement learning, DeepSeekMath-Instruct and DeepSeekMath-RL attain robust results, surpassing 50\% accuracy on the competitive MATH benchmark—a milestone for open-source solutions. In contrast, pre-existing mathematical corpora—typically dominated by English—offer marginal or even negative transfer when it comes to Chinese mathematical reasoning. \paragraph{Natural Language Understanding, Reasoning, and Code}

To evaluate broader language and reasoning abilities, we test the models on MMLU \citep{mmlu} (natural language understanding), BBH \citep{bbh} (reasoning), and examine programming proficiency using HumanEval \citep{codex} and MBPP \citep{mbpp}. For each test question, models are instructed to compose a Python script—permitted to access libraries like \textit{math} and \textit{sympy}—to perform the necessary calculations, and their answers are determined by executing these scripts. \end{itemize}

\subsubsection{Training Setting}
\label{sec:quality-policy}
We fine-tune a general pre-trained language model comprising 1.3B parameters—the same architecture as the DeepSeek LLMs \citep{deepseek-llm}, referred to here as DeepSeek-LLM 1.3B—using mathematical data. We also intend to pursue research directions detailed in Section \ref{sec:effec_rl} to foster more efficient reinforcement learning in large language models. Moreover, we propose Group Relative Policy Optimization (GRPO), a reinforcement learning (RL) algorithm inspired by Proximal Policy Optimization (PPO) \citep{schulman2017proximal}, which eliminates the critic network and instead computes the baseline using group-based scoring. \end{itemize}

\paragraph{Results}
The empirical results presented in Table \ref{tab:code-math} and Table \ref{tab:code-math-general-eval} illustrate model performance across varied training configurations. To tackle these issues, we introduce Group Relative Policy Optimization (GRPO) as depicted in Figure \ref{fig:grpo}. Subsequently, the fastText model is leveraged to retrieve math-relevant pages from this deduplicated set. \paragraph{Algorithms} Algorithms utilize the available data and reward feedback to compute gradient updates for modifying model parameters. \subsection{Training and Evaluating DeepSeekMath-RL}

RL is performed on the DeepSeekMath-Instruct 7B model. \subsubsection{Evaluation Results}

\textbf{The DeepSeekMath~Corpus stands out for its exceptional quality, expansive multilingual coverage, and superior scale relative to comparable datasets.}

\begin{itemize}[topsep=0pt]
    \item \textbf{High-quality}:
    We assess the model’s effectiveness across 8 mathematical benchmarks utilizing few-shot chain-of-thought prompting \cite{cot}. \item We establish a comprehensive framework to interpret a range of methods including RFT, DPO, PPO, and GRPO. \section{Math Pre-Training}

\subsection{Data Collection and Decontamination}

This section details the procedure for curating the DeepSeekMath Corpus from the Common Crawl dataset. In this context, GSM8K and MATH with chain-of-thought responses serve as in-domain benchmarks, whereas all other datasets are categorized as out-of-domain. During training, examples are concatenated at random to approach a maximum context window of 4K tokens. The model undergoes training for 500 iterations, using a batch size of 256 and a fixed learning rate of 5e-5. For an in-depth account of the derivation, consult Appendix \ref{app:analysis-rl}. \subsection{Contributions}
This work makes advances in both large-scale mathematical pre-training and the systematic investigation of reinforcement learning techniques. The construction of this dataset leverages a fastText-based classifier \citep{joulin2016fasttext} to extract relevant data from the Common Crawl (CC). Results in Table \ref{tab:base_math_pot_proof} show that DeepSeekMath-Base 7B sets a new benchmark by exceeding the performance of the former leader, Llemma 34B. We hypothesize this could explain why RL here only boosts Maj@K scores. \item \textbf{Natural Language Understanding, Reasoning, and Code}: To comprehensively evaluate models across language understanding, reasoning, and coding skills, we subject DeepSeekMath-Base to the Massive Multitask Language Understanding (MMLU) benchmark \citep{mmlu}, encompassing 57 multiple-choice subjects, and BIG-Bench Hard (BBH) \citep{bbh}, which is comprised of 23 complex tasks emphasizing multi-step reasoning. To address this, we broaden the diversity of our seed corpus by identifying and incorporating additional mathematical web domains, thereby enhancing the training data for fastText. In the context where models may blend natural language reasoning and computational tool utilization, DeepSeekMath-Instruct 7B attains an accuracy close to 60\% on the MATH benchmark, outperforming all other open-source contenders. Specifically, given a question $q$ and a set of $G$ outputs $\{o_1, o_2, \cdots, o_G\}$, a process-level reward model assigns scores to each reasoning step in the outputs. \subsubsection{Why RL Works?} In our study, reinforcement learning is applied to a subset of instruction tuning data, and this process brings notable performance improvements over the instruction tuning baseline. Each reward is normalized against the global mean and standard deviation as $\widetilde{r}_i^{index(j)} = \frac{r_i^{index(j)} - {\rm mean(\mathbf{R})}}{{\rm std(\mathbf{R})}}$. According to Table \ref{tab:code-math}, in two-stage training, pretraining exclusively on code already provides substantial gains in tackling GSM8K and MATH tasks with Python support, while the subsequent math specialization phase produces further improvement. The term $A_t$ signifies the advantage estimate, computed using Generalized Advantage Estimation (GAE) \citep{gae}, which relies on future rewards $\{r_{\ge t}\}$ and a learned value function $V_{\psi}$. Specifically, during preliminary testing, the model failed on questions involving geometric figures such as triangles and ellipses, potentially pointing to limitations or imbalances in the pretraining and fine-tuning dataset selection. \item Our foundational model, DeepSeekMath-Base 7B, performs on par with Minerva 540B \citep{minerva}, suggesting that parameter count alone does not determine proficiency in mathematical reasoning. Any uncollected web pages connected to these annotated URLs are subsequently incorporated into the seed set. From Figure \ref{fig:rl-analysis}, it is evident that Online RFT yields a substantial performance boost over RFT across two evaluation benchmarks. After this initial phase, the classifier is used to retrieve additional positive samples from CC, which are then subjected to human review for greater accuracy. By the fourth collection round, approximately 98\% of the pages had already been gathered in the prior round, prompting us to halt further collection at this stage. The objective is to generate a formal proof in response to an informal statement, its formal restatement, and an informal proof outline. \end{document} Consistent with DeepSeek LLM training methodologies, we employ the AdamW optimizer \citep{adamW} with hyperparameters $\beta_1=0.9$, $\beta_2=0.95$, and $\mathrm{weight\_decay}=0.1$. For evaluation, we measure mathematical performance of our models with and without tool support on four quantitative reasoning benchmarks presented in English and Chinese. \item \textbf{Online Rejection Sampling Fine-tuning (Online RFT)}: Unlike standard RFT, Online RFT leverages an evolving policy, initializing with the SFT model and repeatedly fine-tuning using outputs sampled from the current policy rather than a fixed model. \paragraph{Data Source} The data source serves as the foundation for any training method. According to Table \ref{tab:base_math_pot_proof}, DeepSeekMath-Base 7B delivers robust performance in automating formal proof generation. Subsequent to pre-training, DeepSeekMath-Base undergoes mathematical instruction tuning utilizing datasets containing chain-of-thought reasoning \citep{cot}, program-of-thought formats \citep{pot,pal}, and tool-augmented reasoning processes \citep{tora}. Consequently, our investigation will center on reinforcement learning algorithms that maintain robustness in the presence of noisy reward feedback. \paragraph{Reward Function} The reward function furnishes the principal signal for learning. Importantly, DeepSeekMath-Base outshines alternatives in Chinese assessments, likely due to our approach of supplementing high-quality multilingual pre-training data rather than restricting to English as in previous studies \citep{minerva,llemma}. The RL training dataset consists of approximately 144K chain-of-thought formatted items drawn from GSM8K and MATH-related SFT questions, while other SFT data is omitted to better assess RL’s effects in settings without overlap between training and evaluation tasks. \end{itemize}

\textbf{One-Stage Training}

\begin{itemize}[topsep=0pt]
    \item \textbf{Math Training for 150B Tokens}: Here, DeepSeek-LLM 1.3B undergoes direct training solely on 150B math tokens;
    \item \textbf{Training on a mixture of 400B Code Tokens and 150B Math Tokens}: Previous observations indicate that math training after code-focused pretraining can undermine code-related capabilities. Annotations for solutions are provided in both CoT and tool-integrated reasoning formats. We investigated this claim across model sizes, namely DeepSeek-LLM 1.3B and DeepSeek-Coder-Base-v1.5 7B \citep{deepseek-coder}, using distinct arXiv-derived corpora processed through different pipelines:

\begin{itemize}[topsep=0pt]
    \item \textbf{MathPile} \citep{mathpile}: an 8.9B-token dataset created with targeted cleaning and filtering heuristics, consisting of more than 85\% scientific articles from arXiv;
    \item \textbf{ArXiv-RedPajama} \citep{redpajama}: a corpus containing all available arXiv LaTeX documents, with preambles, commentary, macros, and bibliographies excluded, summing to 28.0B tokens. Consequently, the PPO algorithm necessitates jointly training a value function alongside the policy network. \end{itemize}

Therefore, more comprehensive investigation is needed, which we defer to future research. These findings partially address the ongoing debate regarding code training’s contribution to reasoning skills: our evidence suggests it does, particularly for mathematical reasoning. \begin{itemize}[topsep=0pt]
    \item \textbf{English mathematical datasets}: We provide tool-integrated annotated solutions for GSM8K and MATH problems, and supplement these with a selected portion of MathInstruct \citep{MathInstruct} and the training partition of Lila-OOD \citep{lila}, wherein solutions utilize CoT or PoT approaches. Unlike traditional methods, GRPO eliminates the need for a critic model, opting instead to derive the baseline from group-based scores. Unlike PPO, GRPO dispenses with explicit value function approximation; instead, it employs the mean reward across a collection of sampled outputs—generated for the same question—as the baseline. To construct the training data, 500,000 samples are drawn randomly from OpenWebMath as positive instances, while an equal number of web pages from Common Crawl are sampled as negatives. \subsubsection{Code Training Benefits Mathematical Reasoning}

A widely held but yet-to-be-proven belief posits that exposure to code enhances reasoning ability. Table \ref{tab:corpora_comparison} clearly illustrates that the model trained on DeepSeekMath~Corpus consistently surpasses those trained on alternative datasets. The optimization follows the objective laid out in equation (\ref{eq:GRPO-obj}). Evaluation for DeepSeekMath-RL 7B is performed on the same benchmarks as those used for DeepSeekMath-Instruct 7B. According to Figure \ref{fig:iter-rl}, iterative RL yields considerable gains, with the most pronounced improvements occurring during the initial iteration. For each training example, $64$ outputs are sampled. The policy is then updated to maximize the objective described in equation (\ref{eq:GRPO-obj}). This allows us to examine the extent to which code data offers specific benefits over general-domain data for mathematical reasoning. Inspired by \cite{wang2023math}, we consider process supervision, which provides feedback after each intermediate reasoning step. Thereafter, process supervision sets the advantage of each token $t$ to be the sum of normalized rewards from all subsequent steps: $\hat{A}_{i, t} = \sum_{index(j) \ge t} \widetilde{r}_i^{index(j)}$. Within the large language model (LLM) setting, it is typical for the reward model to assign a score only to the terminal token, posing challenges for value function training that requires accurate estimation at every token. We assess DeepSeekMath-Base 7B on the informal-to-formal proof generation task proposed by \citep{dsp_proof}. In our investigation, the reward function is characterized as either ‘Rule’-based, evaluating answer correctness, or ‘Model’-based, wherein a learned reward model assigns scores to each response. Formally, the gradient of the training objective with respect to the parameter $\theta$ for these methods is expressible as follows:

\begin{equation}
    \nabla_{\theta}\mathcal{J}_{\textcolor{red}{\mathcal{A}}}(\theta) = \mathbb{E}[\underbrace{(q,o) \sim \textcolor{red}{\mathcal{D}}}_{Data \ Source}]\left( \frac{1}{|o|} \sum_{t=1}^{|o|}  \underbrace{GC_{{\mathcal{A}}}(q, o, t, \textcolor{red}{\pi_{{rf}}})}_{Gradient \ Coefficient}  \nabla_{\theta}\log \pi_{\theta}(o_t | q, o_{<t})\right). Additionally, we observe increased generalization capabilities during reinforcement learning, as evidenced by gains in out-of-domain benchmarks. \end{itemize}

\subsection{Training and Evaluating DeepSeekMath-Base 7B}

This section details the development and evaluation of DeepSeekMath-Base 7B, a foundational model characterized by advanced reasoning proficiency, particularly within mathematics. As detailed in Equation \ref{eq:objective}, three major elements emerge: Data Source, Algorithm, and Reward Function. The revised dataset further informs subsequent iterations of classifier training, thereby enhancing its precision. Our methodology first involves segmenting the entire Common Crawl dataset into mutually exclusive domains, each defined as a collection of web pages under a single base URL. Through this iterative process—conducted four times—we amass a total of 35.5M mathematical web pages, yielding 120B tokens. \item \textbf{PPO/GRPO}: PPO/GRPO strategies initialize the policy with the SFT model and subsequently employ reinforcement learning, sampling outputs from the live policy model to further adapt behavior. Unlike the KL penalty form in (\ref

\begin{algorithm}[t]
  \small
  \caption{Iterative Group Relative Policy Optimization}
  \textbf{Input} initial policy model $\pi_{\theta_{\text{init}}}$; reward models $r_\varphi$; task prompts $\mathcal{D}$; 
  hyperparameters $\varepsilon$, $\beta$, $\mu$
  \begin{algorithmic}[1]
    \State Initialize policy model $\pi_\theta \leftarrow \pi_{\theta_{\text{init}}}$
    \For{iteration = 1, \dots, I}
       \State Assign current policy to reference model: $\pi_{ref} \leftarrow \pi_{\theta}$
      \For{step = 1, \dots, M}
      \State Draw a batch $\mathcal{D}_b$ from $\mathcal{D}$
      \State Assign $\pi_{\theta}$ to previous policy: $\pi_{\theta_{old}} \leftarrow \pi_{\theta}$ 
      \State For each $q \in \mathcal{D}_b$, generate $G$ samples $\{o_i\}_{i=1}^G$ using $\pi_{\theta_{old}} (\cdot \mid q) $
      \State For every sampled $o_i$, evaluate rewards $\{r_i\}_{i=1}^{G}$ by applying $r_{\varphi}$ 
      \State Estimate $\hat{A}_{i,t}$ for the $t$-th token of each $o_i$ utilizing group relative advantage computation. The policy model in GRPO is optimized at a 1e-6 learning rate, and we use a KL penalty coefficient of 0.04. To address this, we investigate an iterative approach to RL using GRPO. This innovation dramatically curtails computational requirements during training. Furthermore, during the RL phase, the value function serves as a baseline for reducing the variance in advantage calculations. Except where explicitly stated, our experiments follow the procedures described in Section \ref{sec:quality-policy}. By expanding the seed corpus in this targeted manner, we substantially increase the quantity of positive examples available, which strengthens the fastText model's ability to retrieve mathematical documents in subsequent cycles. \paragraph{Observation about Gradient Coefficient} The mechanism by which reward information modulates the gradient coefficient for parameter updates varies across algorithms. This lead persists even when comparing against significantly larger models (e.g., Qwen 72B) or those explicitly trained with reinforcement learning methods tailored to mathematical tasks (e.g., WizardMath-v1.1 7B). The learning rate follows a multi-stage schedule: after an initial 2,000-step warmup, it ascends to its maximum, then drops to 31.6\% at the 80\% mark of the training, and finally declines to 10.0\% of its peak value at 90\% of the training completion. The training is executed using an open-source implementation\footnote{\url{https://fasttext.cc}} with a vector dimension of 256, a learning rate of 0.1, a maximum n-gram length of 3, a minimum word frequency threshold of 3, and 3 epochs. This held true across various tasks ranging in complexity, encompassing quantitative reasoning benchmarks such as GSM8K and MATH (see Table \ref{tab:arxiv-cot}), multiple-choice tests like MMLU-STEM (refer to Table \ref{tab:arxiv-cot}), as well as formal math problems (miniF2F, Table \ref{tab:arxiv_atp}). In our experiments with Proof-Pile-2, we adhere to the 2:4:1 ratio of arXiv to web to code tokens as suggested in \cite{llemma}. Together, these cover a comprehensive range of mathematical subjects from elementary through collegiate levels. The dataset composition is distributed as follows: 56\% originates from the DeepSeekMath Corpus, 4\% from AlgebraicStack, 10\% from arXiv publications, 20\% consists of Github code, and the remaining 10\% is derived from English and Chinese text in the Common Crawl dataset. In total, our dataset consists of 776K training samples. Moving forward, we aim to further refine our data selection and curation processes to assemble an even higher-quality pretraining dataset. We anticipate that alignment strategies exemplified by \textbf{WEAK-TO-STRONG} \citep{burns2023weak} have the potential to transform the core mechanisms of learning algorithms. DeepSeekMath~7B demonstrates strong results, attaining a 51.7\% score on the highly challenging MATH benchmark—achieved without the assistance of supplementary toolkits or voting strategies—narrowing the gap with state-of-the-art models such as Gemini-Ultra and GPT-4. This work confines itself to questions from the instruction tuning phase and relies on basic nucleus sampling for output generation. \section{Conclusion, Limitation, and Future Work} 

This work introduces DeepSeekMath, a model that surpasses all other open-source systems on the challenging MATH benchmark and nearly matches the proficiency of proprietary counterparts. \subsubsection{Process Supervision RL with GRPO} While outcome supervision relies on terminal rewards, this approach may lack granularity for supervising complex reasoning tasks. This set includes not only the popular general-purpose Mistral 7B \citep{mistral} but also Llemma 34B \citep{llemma}, a recently introduced model further trained on the Proof-Pile-2 corpus \citep{llemma} for mathematical tasks. Our English dataset reflects a variety of mathematical domains, such as algebra, probability, number theory, calculus, and geometry. We propose that our approach to mathematical data curation can serve as a foundation for further scholarly exploration, as considerable potential for progress remains. This design choice leads to a substantial reduction in computational demands in comparison to Proximal Policy Optimization (PPO). In subsequent studies, we aim to deploy our RL approach with out-of-distribution questions, integrating \textbf{advanced sampling (decoding) strategies} such as those leveraging tree-search algorithms \citep{yao2023tree}. \EndFor 
  \end{algorithmic}
  \textbf{Output} $\pi_\theta$
  \label{alg:iter-grpo}
\end{algorithm}

\subsubsection{Outcome Supervision RL with GRPO} In this framework, for any given question $q$, the method samples a collection $\{o_1, o_2, \cdots, o_G\}$ of responses from the previous policy model $\pi_{\theta_{old}}$. As reported in Table \ref{tab:corpora_comparison}, leveraging DeepSeekMath~Corpus during training leads to enhanced mathematical reasoning in both English and Chinese. \subsection{Lessons Learnt in Pre-Training}

We begin by discussing our pre-training observations. \paragraph{Observation about Data Source} We categorize the origin of training data into online and offline sampling. Utilizing only a selected portion of English instruction tuning data, GRPO achieves considerable improvements over the strong DeepSeekMath-Instruct model, including significant gains on both in-domain tasks (GSM8K: 82.9\% $\rightarrow$ 88.2\%, MATH: 46.8\% $\rightarrow$ 51.7\%) and out-of-domain mathematical challenges (e.g., CMATH: 84.6\% $\rightarrow$ 88.8\%) during RL training. Notably, this strategy can be generalized to other specialized domains such as programming. Nevertheless, there is a paucity of in-depth studies on how these resources actually affect models’ capabilities in mathematical reasoning. The rewards are standardized by removing the mean of the group and dividing by its standard deviation. \section{Reinforcement Learning}

\subsection{Group Relative Policy Optimization} Reinforcement learning (RL) methods have demonstrated substantial efficacy in enhancing the mathematical reasoning proficiency of LLMs post-Supervised Fine-Tuning (SFT) \citep{wang2023math,wizardmath}. To mitigate the risk of contaminating downstream benchmarks, we adopt the filtering protocols described in \cite{deepseek-coder}: we exclude any web pages featuring questions or answers from well-known mathematical evaluation sets, including English benchmarks like GSM8K \citep{gsm8k} and MATH \citep{MATH}, as well as Chinese datasets such as CMATH \citep{wei2023cmath} and AGIEval \citep{agieval}. Results indicate that the inclusion of arXiv papers, in isolation, does not bolster mathematical reasoning: models trained solely on these corpora showed no measurable improvement and sometimes even exhibited worsened performance on a diverse suite of mathematical evaluations. We identify three pivotal research avenues for reward models: 1) \textbf{Strategies for improving reward model generalization.} Effective reward models need to generalize across out-of-distribution queries and novel decoding outputs; failing to do so may cause reinforcement learning to merely reinforce the status quo of LLM outputs without substantive advancement; 2) \textbf{Methods for representing reward model uncertainty.} Accurately capturing uncertainty may offer a crucial interface between unreliable reward functions and weak-to-strong learning approaches; 3) \textbf{Approaches to constructing high-quality process reward models efficiently} so as to deliver granular feedback signals throughout reasoning steps \citep{lightman2023let,wang2023math}. Additionally, we provide a cohesive framework for interpreting multiple training strategies and outline several promising research avenues to develop more advanced reinforcement learning methods. Additionally, HumanEval \citep{codex} and MBPP \citep{mbpp} are used as standard benchmarks for code evaluation. Here, we present DeepSeekMath~7B, which builds upon DeepSeek-Coder-Base-v1.5 7B by conducting additional pre-training on 120B mathematics-specific tokens collected from Common Crawl, complemented by data from natural language and programming code sources. Secondly, the proposal of Group Relative Policy Optimization (GRPO), an extension of Proximal Policy Optimization (PPO), which strengthens the model’s mathematical reasoning skills and, at the same time, leads to more efficient memory utilization during PPO. The reward model is further refined using a replay buffer that mixes 10\% of prior data with the new samples. \item While utilizing arXiv documents in training is a prevalent strategy—especially in mathematics-related research—our results reveal it yields no discernible performance gains across the mathematical benchmarks considered in this study. Initially, OpenWebMath \citep{openwebmath}, a reputable source of mathematical web content, is adopted as the seed corpus. Online sampling refers to data collected directly from the on-policy exploration of the live training model, whereas offline sampling relies on data generated by the pretrained SFT model. Additionally, we find that math-specific training not only advances mathematical competencies, but also yields gains on broader reasoning benchmarks, such as MMLU \citep{mmlu} and BBH \citep{bbh}, demonstrating improvements in both specialized and general reasoning abilities. Moreover, thanks to its multilingual content, the DeepSeekMath Corpus contributes to notable gains in Chinese mathematical benchmark tasks \citep{wei2023cmath,agieval}. Central to this advancement is the development of the DeepSeekMath Corpus—a comprehensive, high-quality pre-training dataset consisting of 120B math tokens. Specifically, given a question $q$, GRPO draws a set of outputs $\{o_1, o_2, \cdots, o_G\}$ from the previous policy $\pi_{\theta_{old}}$ and seeks to optimize the policy according to the following objective:

\begin{equation}
\footnotesize
\begin{split}
    \mathcal{J}_{GRPO}(\theta) &= \mathbb{E}{[q \sim P(Q), \{o_i\}_{i=1}^G \sim \pi_{\theta_{old}}(O|q)]}  \\
    & \frac{1}{G}\sum_{i=1}^G\frac{1}{|o_i|} \sum_{t=1}^{|o_i|} \left\{ \min \left[ \frac{\pi_\theta(o_{i,t} | q, o_{i,<t})}{\pi_{\theta_{old}}(o_{i,t} | q, o_{i,<t})} \hat{A}_{i,t}, \text{clip} \left( \frac{\pi_\theta(o_{i,t} | q, o_{i,<t})}{\pi_{\theta_{old}}(o_{i,t} | q, o_{i,<t})}, 1 - \varepsilon, 1 + \varepsilon \right)  \hat{A}_{i,t} \right] - \beta \mathbb{D}_{KL}\left[\pi_{\theta} || \pi_{ref}\right]\right\} ,
\end{split}
\label{eq:GRPO-obj}
\end{equation}

Here, $\varepsilon$ and $\beta$ are treated as hyper-parameters, and the term $\hat{A}_{i,t}$ signifies the advantage, which is derived using only the relative rewards among outputs sampled within the same group—a process that will be further detailed in the subsequent sections. We hypothesize that this tradeoff arises from the constrained model capacity of DeepSeek-LLM 1.3B, which may not suffice for the concurrent learning of both code and mathematics at once. \subsubsection{How to Achieve More Effective RL?}
\label{sec:effec_rl}
Our findings establish that RL is quite effective for mathematical reasoning tasks. Additionally, these systems have emerged as valuable tools for assisting humans with complex mathematical challenges \citep{tao}. \item \textbf{Direct Preference Optimization (DPO)}: DPO enhances the SFT model through supplementary training on paired outputs produced by the SFT model, optimizing according to a pairwise DPO loss criterion. \end{abstract}

\section{Introduction}

Large language models (LLMs) have transformed how artificial intelligence tackles mathematical reasoning, leading to marked progress in both quantitative reasoning benchmarks \citep{MATH} and those focused on geometric reasoning \citep{trinh2024solving}. The training data for the reward model originates from rule-based evaluation. \subsubsection{ArXiv Papers Appear to Offer Limited Gains for Mathematical Reasoning}

ArXiv papers frequently constitute a substantial portion of mathematical pre-training datasets \citep{minerva,gpt-f,llemma,mathpile}. The group-based method used by GRPO for advantage estimation is highly compatible with reward models that are themselves trained using pairwise comparisons on a single question. \item Drawing on insights from this unified framework, we investigate what drives the effectiveness of reinforcement learning and highlight several promising research avenues aimed at advancing reinforcement learning for large language models (LLMs). \subsection{Validating the Quality of the DeepSeekMath~Corpus}

To evaluate the merits of the DeepSeekMath~Corpus, we conduct pre-training studies that compare it against other recent large-scale mathematical training datasets:

\begin{itemize}[topsep=0pt]
    \item \textbf{MathPile} \citep{mathpile}: a composite dataset totaling 8.9B tokens, aggregated from sources such as textbooks, Wikipedia, ProofWiki, CommonCrawl, StackExchange, and arXiv—with over 85\% of the content derived from arXiv;
    \item \textbf{OpenWebMath} \citep{openwebmath}: a dataset filtered from CommonCrawl to retain only mathematical web pages, comprising 13.6B tokens;
    \item \textbf{Proof-Pile-2} \citep{llemma}: a curated corpus that integrates OpenWebMath, AlgebraicStack (with 10.3B tokens of mathematical code), and arXiv articles (amounting to 28.0B tokens). However, a one-stage approach with combined code and math data impedes the model’s mathematical reasoning abilities absent tool usage. For context, any reference to the DeepSeekMath Corpus in this discussion refers to the 89-billion-token dataset generated during the second round of data collection. In two-stage setups, the initial code phase contributes measurable improvement on its own, and also enhances the efficacy of later math-focused training, culminating in superior results. According to Table \ref{tab:sft_rl_math}, when tools are not available during evaluation, DeepSeekMath-Instruct 7B delivers robust stepwise reasoning abilities. (2) Notably, DeepSeekMath-RL 7B is exclusively fine-tuned with chain-of-thought-format instruction data from GSM8K and MATH, originating from DeepSeekMath-Instruct 7B. We further introduce a unified framework that contextualizes various approaches—such as Rejection Sampling Fine-Tuning (RFT) \citep{yuan2023scaling}, Direct Preference Optimization (DPO) \citep{dpo}, PPO, and GRPO—under a spectrum of direct or simplified RL techniques. \end{itemize}

The elements of these approaches are outlined in Table \ref{tab:data_policy}. As depicted in Algorithm \ref{alg:iter-grpo}, this method involves regenerating training datasets for the reward model by sampling from the latest policy model. Figure \ref{fig:maj-pass} shows that RL leads to marked gains in Maj@K while leaving Pass@K largely unaffected. Furthermore, we propose Group Relative Policy Optimization (GRPO), a tailored modification of Proximal Policy Optimization (PPO), which significantly enhances mathematical reasoning while reducing memory usage. Through comprehensive ablation studies, we determine that online web sources are particularly fruitful for obtaining rich mathematical data, whereas data from arXiv does not contribute as substantially as previously anticipated. Each problem is accompanied by its solution, formatted as chain-of-thought (CoT) \citep{cot}, program-of-thought (PoT) \citep{pot,pal}, or tool-integrated reasoning, as per \citep{tora}. For the initial training, the classifier uses samples from OpenWebMath \citep{openwebmath} as positive instances, with a diverse array of unrelated web content serving as negative cases. Construction of reward model training data follows the method outlined by \citep{wang2023math}. This paper presents DeepSeekMath, a language model tailored specifically for mathematical domains, which demonstrates a substantial improvement over open-source alternatives and closely matches GPT-4’s results on academic evaluations. The peak learning rate is set to 5.3e-4, with training performed using a batch size of 4M tokens and a context window of 4K tokens. In particular, during the initial training phase, both methods perform similarly; however, as training progresses, Online RFT distinctly surpasses RFT, thereby demonstrating the effectiveness of online data collection. As training advances, samples drawn from the actor diverge more notably, amplifying the benefits of adaptive, real-time data acquisition. Nevertheless, it is infeasible to guarantee that the reward signals are perpetually accurate, particularly when addressing highly intricate tasks. Furthermore, code-oriented pretraining strengthens mathematical reasoning skills even when the use of external tools is not involved. DeepSeekMath-Instruct is competitive with top-tier Chinese proprietary models, including GLM-4 and Baichuan-3, though it still lags behind leading systems like GPT-4 and Gemini Ultra on MATH. In outcome supervision, the normalized reward—calculated at the end of every output $o_i$—is used as the advantage for every token in the output: $\hat{A}_{i, t} = \widetilde{r}_i = \frac{r_i- {\rm mean}(\mathbf{r})}{{\rm std}(\mathbf{r})}$. Figure \ref{fig:corpora_comparisons} shows that models trained with this corpus, such as DeepSeek-LLM 1.3B, achieve sharper gains in learning as well as sustained performance improvements. \subsubsection{Iterative RL with GRPO} During reinforcement learning, the existing reward model might no longer provide adequate supervision as the policy model evolves. The model is built upon DeepSeek-Coder-v1.5 7B as its base, receiving continual pretraining over 500B tokens, including an extensive 120B-token subset specifically containing mathematical content extracted from Common Crawl. Moreover, \textbf{efficient inference techniques} \citep{xia-etal-2023-speculative,leviathan2023fast,kwon2023efficient,xia2024unlocking}, which critically influence how effectively the policy model explores, are anticipated to play a pivotal role. After this, the policy model becomes the new reference model, and is further updated based on feedback from the updated reward model. \item \textbf{Rejection Sampling Fine-tuning (RFT)}: RFT builds on the SFT model by further training it using output samples generated in response to SFT prompts, keeping only those outputs deemed correct. Each response is evaluated with a reward model, generating the reward set $\mathbf{r}=\{r_1, r_2, \cdots, r_G\}$. For each distinct mathematical dataset, the model is trained independently for a total of 150B tokens. Below, we discuss prominent approaches as instantiated within this framework:

\begin{itemize}
    \item  \textbf{Supervised Fine-tuning (SFT)}: In SFT, the pretrained model undergoes additional training using manually curated SFT datasets. For benchmark text snippets shorter than 10 grams but longer than 3 grams, we implement exact string matching to filter out the contaminated web material. Yet, the highest-performing systems such as GPT-4 \citep{gpt4} and Gemini-Ultra \citep{gemini} remain proprietary, with the present open-source alternatives falling substantially short in their comparative capabilities. \item \textbf{Formal Mathematics}: DeepSeekMath-Base is assessed on the informal-to-formal theorem proving benchmark from \citep{dsp_proof}, utilizing miniF2F \citep{minif2f} with Isabelle \citep{isabelle} serving as the selected proof assistant. \For{GRPO iteration = 1, \dots, $\mu$}
        \State Improve the policy model $\pi_{\theta}$ by optimizing the GRPO criterion (Equation \ref{eq:GC-GRPO})
      \EndFor
    \EndFor 
    \State Continuously update $r_\varphi$ by training with a replay buffer. \paragraph{Mathematical Problem Solving with Step-by-Step Reasoning}

DeepSeekMath-Base's effectiveness in mathematical problem solving is assessed via few-shot chain-of-thought prompting \citep{cot}, spanning eight benchmarks in both English and Chinese. \end{itemize}

\subsection{Training and Evaluating DeepSeekMath-Instruct 7B}

This section outlines the development of DeepSeekMath-Instruct 7B, which is further trained for mathematical instruction tasks starting from DeepSeekMath-Base. In reinforcement learning settings, this role is typically filled by a neural reward model. Illustrated in Figure \ref{fig:data_collect}, an iterative workflow is employed to systematically extract a large-scale mathematical text corpus, beginning with a seed dataset (i.e., a curated set of high-quality math resources). \item Insights from our experimental math training indicate that conducting code pre-training before mathematical pre-training notably enhances models' mathematical problem-solving capacities, regardless of the presence of tool augmentation. For the initial phase, the dataset is restricted to the top 40B tokens. To counteract excessive optimization of the reward model, a commonly adopted strategy involves incorporating a per-token KL penalty—relative to a reference model—directly into the reward at every token, following \citep{ouyang2022training}:

\begin{equation}
    r_{t} = r_\varphi(q, o_{\le t}) - \beta \log\frac{\pi_{\theta}(o_{t}|q, o_{<t})}{\pi_{ref}(o_{t}|q, o_{<t})},
\label{eq:PPO-reward}
\end{equation}

In this expression, $r_\varphi$ stands for the reward model, and $\pi_{ref}$ represents the reference model (typically the initial SFT model), while $\beta$ acts as the weighting coefficient for the KL penalty. Our analysis yields several insights: (1) With chain-of-thought reasoning, DeepSeekMath-RL 7B achieves accuracy rates of 88.2\% on GSM8K and 51.7\% on MATH. Furthermore, Figure \ref{fig:corpora_comparisons} demonstrates that after 50B training tokens (representing a complete pass over Proof-Pile-2), DeepSeekMath-trained models exhibit superior performance, evidencing the higher overall quality of the DeepSeekMath~Corpus. \paragraph{Formal Mathematics}

Automating formal proof construction offers significant benefits in both verification accuracy and productivity, which has driven increased research in this area. \end{itemize}

To assess their effect, we independently pre-trained DeepSeek-LLM 1.3B for 150B tokens and DeepSeek-Coder-Base-v1.5 7B for 40B tokens using each arXiv dataset. When employing self-consistency with 64 sampled completions, DeepSeekMath~7B reaches 60.9\% accuracy on MATH. Even with such a limited fine-tuning dataset, it delivers superior results compared to DeepSeekMath-Instruct 7B in every metric evaluated, illustrating the substantial gains provided by reinforcement learning. Our evaluation covers both the ability to produce comprehensive textual solutions unaided by external tools and the capacity to answer via Python computation. To investigate the impact of code training on mathematical reasoning capabilities, we implemented and compared two distinct training paradigms: a two-stage training approach and a single-stage training approach. DeepSeekMath-Base is initialized from DeepSeek-Coder-Base-v1.5 7B \citep{deepseek-coder}, given that models pre-trained on code exhibit superior performance relative to those initialized from general language models. In accordance with \cite{dsp_proof}, our protocol has the model draft proof sketches, after which Sledgehammer \citep{sledgehammer}, an established automated prover, completes any unresolved components. The selected benchmarks represent a spectrum of tasks, including quantitative reasoning—such as GSM8K \citep{gsm8k}, MATH \citep{MATH}, and CMATH \citep{wei2023cmath}—as well as multiple-choice assessments like MMLU-STEM \citep{mmlu} and Gaokao-MathQA \citep{agieval}. It is important to emphasize, however, that these results are not exhaustive and should be interpreted cautiously. This yields DeepSeekMath-Instruct 7B, a model that outperforms all existing 7B instruction-tuned counterparts and attains performance comparable to 70B open-source instruction-tuned models. Domains where over 10\% of the pages were included are deemed math-centric (e.g., \url{mathoverflow.net}). \section{Discussion}
This section details our key observations stemming from both pre-training and reinforcement learning studies. The model's initial parameters are sourced from DeepSeek-Coder-Base-v1.5 7B \citep{deepseek-coder}, and it undergoes training on a corpus comprising 500B tokens. To calibrate the preserved dataset size, pre-training trials are conducted on collections with 40B, 80B, 120B, and 160B top-ranked tokens. The model’s proficiency in mathematical reasoning is largely credited to two main elements: Firstly, the robust exploitation of freely available internet data, achieved through a carefully developed data selection pipeline. \textbf{Math Pre-Training at Scale}

\begin{itemize}[topsep=0pt]
    \item We demonstrate that the openly available Common Crawl dataset holds a substantial amount of mathematical knowledge. Finally, we elucidate the reasons behind the efficacy of RL in enhancing instruction-tuned model performance, and outline promising future directions for advancing RL methods within this framework. Additionally, the results reveal that GRPO+PS achieves better outcomes than GRPO+OS, underscoring the value of adopting step-aware, fine-grained gradient coefficient adjustments. To further elucidate the mechanism behind RL's effectiveness, we assess both Pass@K and Maj@K accuracies for Instruct and RL-trained models across two evaluation datasets. Specifically, in RL settings, the data source refers to unlabeled questions accompanied by responses sampled from the policy model. \end{itemize}

\textbf{Exploration and Analysis of Reinforcement Learning}

\begin{itemize}[topsep=0pt]
    \item We propose Group Relative Policy Optimization (GRPO), a novel and resource-efficient reinforcement learning approach. On English tasks, DeepSeekMath-Base demonstrates performance on par with the proprietary Minerva 540B \citep{minerva}, while consistently outperforming all open-source base models (such as Mistral 7B \citep{mistral} and Llemma-34B \citep{llemma}), regardless of prior mathematical pre-training, often by considerable margins. This dataset considerably surpasses prior collections—nearly 7 times larger than the Minerva web data \citep{minerva} and about 9 times greater than the recently introduced OpenWebMath \citep{openwebmath}. According to Equation \ref{eq:objective}, nearly all current approaches essentially place full \textbf{TRUST} in the signals provided by the reward function to either boost or diminish the conditional likelihood of specific tokens. With regard to training methodology, we largely adhere to the protocol described in Section \ref{sec:quality-policy}, though we adjust the peak learning rate to 4.2e-4 and employ a batch size of 10 million tokens. Following the initial round of data gathering, a significant number of mathematical web pages remain absent from the corpus, largely because the fastText model had originally been trained on a set of positive samples lacking sufficient variety. On other evaluated datasets, it remains on par with DeepSeek-LLM-Chat 67B, the previous leading model in its class, despite being only one-tenth its size. Empirical evaluation demonstrates the corpus’s quality, with DeepSeekMath-Base 7B achieving 64.2\% accuracy on GSM8K \citep{gsm8k} and 36.2\% on the competition-grade MATH benchmark \citep{MATH}, surpassing the performance of Minerva 540B \citep{minerva}. Specifically, this work does not address:

\begin{itemize}[topsep=0pt]
    \item How arXiv data may impact specialized math tasks beyond those studied here, such as translating formal mathematical statements or proofs into informal language;
    \item The potential benefits of arXiv data when mixed with other categories of training data;
    \item Whether advantages might become apparent at much larger model scales. We also investigate iterative RL by running two cycles in our experimental setup. We additionally introduce a unified perspective that encapsulates prominent training strategies, categorizing them as either direct or approximate forms of RL. The variables $q$ and $o$ refer to questions and outputs, which are sampled both from a dataset of questions and from the old policy $\pi_{\theta_{old}}$. Meanwhile, models trained on substantially smaller corpora quickly encounter data repetition and thus see their accuracy plateau at an earlier stage. Furthermore, it even outperforms Minerva 540B \citep{minerva}, a proprietary model built on PaLM \citep{palm} and further specialized on mathematical texts, despite the latter’s much larger size (77 times greater). To supplement this, we also present an overview of the model’s broader abilities, such as understanding and reasoning in natural language, as well as programming tasks. The evaluation uses miniF2F \citep{minif2f}, a benchmark focusing on Olympiad-grade formal mathematics, and employs few-shot prompts to generate formal Isabelle proofs. Comprehensive experiments, including comparisons between online and offline training, outcome versus process supervision, and single-turn versus iterative RL, are conducted to uncover the fundamental aspects of this unified paradigm. Effective results can be achieved by comparatively smaller models if the pre-training data is of superior quality. Conversely, Online RFT does not integrate this mechanism; it lacks explicit penalty for erroneous outputs and rewards all correct responses identically, irrespective of relative quality. Of particular note, DeepSeekMath-Base exceeds previous open-source base models on the MATH dataset—considered to be at competition level—by a margin exceeding 10\% in absolute terms. The parameter $\varepsilon$ is a hyper-parameter associated with clipping in PPO, which aids in achieving training stability. The need to train a separate value function in PPO, which is generally as large as the policy model itself, results in increased memory and computational demands. Table \ref{tab:base_math_cot} demonstrates that DeepSeekMath-Base 7B achieves the highest scores across all eight benchmarks compared to other open-source base models. Collectively, these results show that DeepSeekMath-Base 7B significantly outstrips the generalist Mistral 7B \citep{mistral} on all three evaluated benchmarks involving reasoning and coding. Moreover, whereas PPO adds a KL penalty term inside the reward function, GRPO incorporates KL regularization directly in the objective as a penalty between the current policy and the reference policy, which keeps the calculation of $\hat{A}_{i,t}$ straightforward. After each stage of exploration, the policy model receives a single update. This pattern suggests that reinforcement learning mainly strengthens the robustness of the model’s output distribution, meaning \textbf{the benefit comes from promoting correct responses within the TopK outputs rather than enhancing the model’s core skills.} Analogously, \citep{wang2023making} reported a \textbf{misalignment problem} in SFT model reasoning tasks, and found that aligning preferences could improve reasoning results, as demonstrated in works such as \citep{yuan2023rrhf,song2023preference,wang2023making}. \item \textbf{Large-scale}:
    The size of DeepSeekMath~Corpus far exceeds that of previously released mathematical corpora. We perform thorough experimental studies—such as comparing online versus offline learning, contrasting supervision over outcomes and processes, and analyzing single-turn versus iterative reinforcement learning—to systematically dissect the core factors in this unified approach. \section{Supervised Fine-Tuning}

\subsection{SFT Data Curation}

We assembled a comprehensive dataset for mathematical instruction tuning, encompassing problems in both English and Chinese, sourced from a broad spectrum of mathematical disciplines and spanning different difficulty levels. \textbf{Two-Stage Training}

\begin{itemize}[topsep=0pt]
    \item \textbf{Code Training for 400B Tokens $\rightarrow$ Math Training for 150B Tokens}: The DeepSeek-LLM 1.3B model is first exposed to 400B tokens of code data, and subsequently trained on 150B tokens of math content;
    \item \textbf{General Training for 400B Tokens $\rightarrow$ Math Training for 150B Tokens}: For comparison, we replace code tokens with general-domain tokens, sampled from a comprehensive general corpus assembled by DeepSeek-AI, during the initial training stage. \subsubsection{From PPO to GRPO} Proximal Policy Optimization (PPO) \citep{schulman2017proximal}, an actor-critic RL framework, has become standard practice in the RL-based fine-tuning phase for LLMs \citep{ouyang2022training}. Table \ref{tab:sft_rl_math} presents a comparison of open- and closed-source models employing both chain-of-thought and tool-augmented reasoning approaches across English and Chinese benchmark datasets.